\chapter{Revisão da Literatura}

%TODO: ESCOLHER MELHORES RELACIONADOS

%https://www.tandfonline.com/doi/full/10.1080/1369118X.2020.1776370#abstract
%https://www.mdpi.com/2079-9292/13/3/648 - ESPECIFICAMENTE CAPITULO 4 E 5
%https://citeseerx.ist.psu.edu/document?repid=rep1&type=pdf&doi=aad01b2a642711ef0b4d7d89d8d50fc268a222ce - ASSUNTO GERAL SOBRE NLP
%https://ijlllc.org/uploads2023/LLLC_02_049.pdf - TRADUÇÃO DE INGLES PARA IGBO DE ARQUIVOS LEGAIS
%https://citeseerx.ist.psu.edu/document?repid=rep1&type=pdf&doi=61d09d0d3aad012727e929c3cb539ddd340c9d3b USO DE NLP - MEDICINA
%https://link.springer.com/book/10.1007/978-3-031-02167-1 NLP EM REDE SOCIAL - CATEGORIZAÇÃO/MARKETING ETC (NAO MUITO A VER)
%https://arxiv.org/abs/2103.11878 - AVALIAÇÃO/METRICA PARA AVALIAR TRADUÇÕES DE MÁQUINA AGNOSTICAS A CONTEXTO


\section{Arquiteturas de sistemas de tradução automática}

%TODO: MT OU TM?
%AS PAGINAS DO ARTIGO DO CHÉRAGUI ESTÃO TODAS NA CASA DOS 160, APESAR DE QUE NO PDF ONLINE SÓ TEM UMAS 30 PÁGINAS, EU ESTOU ASSUMINDO QUE ELA FAÇA PARTE DE ALGUM ARTIGO MAIOR MAS EU HONESTAMENTE NÃO ENTENDI...
A Tradução de Máquina (MT), que consiste na conversão de uma língua de origem para uma língua alvo de forma mecânica, evoluiu significativamente ao longo dos anos. Entre as principais arquiteturas de sistemas de tradução automática discutidas na literatura, destacam-se: Abordagens diretas, baseada em transferência e Interlingua \cite{cheragui}.
%(CHÉRAGUI, 2012).

%\begin{itemize} 

\subsection{\textbf{Abordagem direta}}
A abordagem direta realiza a tradução diretamente da língua de origem para a língua alvo. Nessa arquitetura, o sistema analisa o vocabulário da língua de origem para resolver ambiguidades e identificar a ordem correta das palavras na língua alvo. Um dicionário bilíngue e programas de análise léxica e morfológica são fundamentais para esta abordagem \cite[p. 163]{cheragui}
%(CHÉRAGUI, 2012, p. 163).

\subsection{\textbf{Abordagem baseada em transferência}}
A abordagem de transferência utiliza três etapas principais para realizar a tradução. Primeiramente, o texto na língua de origem é convertido em uma representação intermediária; em seguida, essa representação é transformada para uma estrutura equivalente na língua alvo. Por fim, a tradução final é gerada a partir dessa estrutura \cite{cheragui}. A representação intermediária pode variar de uma árvore sintática (\figurename~\ref{fig:arvore-sintatica}) até uma representação mais semântica, dependendo da complexidade e dos requisitos linguísticos.

%https://alexandrehefren.wordpress.com/2010/03/14/determinismo-e-gramatica-sintagmatica-gs-parte-1/
    \begin{center}
\begin{figure}

    \Tree [.F 
        [.SN 
            [.Q O ]
            [.N gato ]
        ]
        [.SV 
            [.V come ]
            [.N peixe ]
        ]
    ]
    
    \caption{Árvore sintática} 
    \label{fig:arvore-sintatica}
    \begin{tabular}{r@{: }l r@{: }l}
    $F$ & Frase & $SN$ & Sintagma Nominal\\
    $SV$ & Sintagma Verbal & $V$ & Verbo\\
    $Q$ & Quantificador & $N$ & Nome 
    \end{tabular}
    

\end{figure}
    \end{center}

\subsection{\textbf{Abordagem Interlingua}}
Considerada ideal para sistemas multilíngues, a abordagem Interlingua realiza a análise do texto de origem, extraindo seu conteúdo semântico e o representando em uma linguagem neutra ("Interlingua") e independente das línguas de origem e destino . Posteriormente, o conteúdo é traduzido para a língua alvo a partir dessa representação neutra, permitindo o uso de uma mesma análise para múltiplas línguas de origem e destinos variados \cite[p. 164]{cheragui}
%(CHÉRAGUI, 2012, p. 164).



%\end{itemize}


% \subsubsection{}
%\paragraph{}

\section{A Importância da tradução de máquina para o contexto jurídico}

Como discutido anteriormente, a estrutura sintática de uma frase e a maneira como diferentes elementos se conectam desempenham um papel fundamental na tradução automática. No entanto, a precisão na tradução se torna ainda mais complexa em contextos que exigem uma interpretação rigorosa e sensível a nuances, como o contexto jurídico. 

\subsection{\textbf{Percepção do uso de MT e riscos no contexto jurídico}}
De acordo com \citeonline{VHS}, a percepção sobre os riscos associados ao uso de MT no setor jurídico varia significativamente \cite[p. 1522]{VHS}.
%(Vieira, O’Hagan e O’Sullivan, 2021, p. 1522).
%\citeonline{VHS}
%CITAÇÃO DESSES 3 AUTORES EM CONJUNTO NÃO ESTÁ FUNCIONANDO DIREITO,

Por exemplo, o Tribunal Estadual do Novo México, nos EUA, se destaca por ter estabelecido diretrizes claras em relação ao uso de MT, por conta de sua frequente circunstância de apontar jurados não fluentes em inglês. Segundo essas diretrizes, traduções feitas por MT que não forem editadas por um humano não devem ser utilizadas em documentos formais, como procedimentos judiciais ou provas \cite[p. 323]{chavez}
%(CHÁVEZ, 2008, p. 323).

%CHAVEZ, 2008)

Evidenciado ainda por \citeonline{VHS}, há em contrapartida, evidências que sugerem uma alta confiança para com a tradução automática, que demonstra para os autores, uma demonstração de falta de conscientização entre profissionais sobre os riscos de MT \cite[p. 1522]{VHS}.
%(Vieira, O’Hagan e O’Sullivan, 2021, p. 1522)
%\cite{Vasquez2019}, seção \textit{Analysis I. Guilty Plea B. Translation}
No caso jurídico de Ernesto Vasquez, em Dallas - Texas, um prisioneiro federal de língua espanhola tentou retirar uma confissão alegando que não compreendeu adequadamente a situação devido a traduções deficientes oferecidas por seu advogado. Curiosamente, o novo advogado que defendeu a retirada da confissão criticou o trabalho do colega anterior, mencionando o Google Translate como uma fonte viável: \textit{“Se eu não conheço a palavra, meritíssimo, eu procuro a tradução. Está disponível gratuitamente no Google. Existe um aplicativo chamado Google Translate”}\footnote{Vasquez vs USA. 2019. Disponível em: \url{https://scholar.google.co.uk/scholar_case?case=4394896296660864715&hl=pt-BR&as_sdt=0}\label{Vasquez}}. Em outro exemplo, o próprio tribunal recorreu ao uso de MT: \textit{“Como os Requerentes não forneceram tradução para nenhum dos documentos em polonês submetidos em apoio à sua petição, o Tribunal usou o serviço gratuito Google Translate para confirmar certas declarações”}\footnote{
Super Express USA Publishing Corp. vs Spring Publishing Corp. 2017. Disponível em: \url{https://casetext.com/case/super-express-us-publg-corp-v-spring-publg-corp}\label{SuperExpress2017}}.
% \cite{SuperExpress2017

Esses exemplos indicam que, apesar dos avanços tecnológicos, a MT ainda apresenta riscos significativos para o setor jurídico, especialmente em situações onde a precisão e a compreensão contextual são fundamentais.

\subsection{\textbf{Aplicações práticas já aceitas de MT no contexto jurídico}}

Nem todo uso de MT é mal-visto. Segundo \citeonline{VHS}, em fases de pré-julgamento, como a etapa de descoberta — em que as partes trocam evidências e informações legais —, o risco de usar MT é considerado menor. Nesse contexto, a MT pode atuar como uma ferramenta inicial de triagem para revisão rápida de informações eletronicamente disponíveis, facilitando a identificação de documentos relevantes para análise subsequente por tradutores profissionais \cites[p. 45]{foster}[p. 19]{nelson}
%(FOSTER; NORTHROP, 2011, p. 45; NELSON; SIMEK, 2018, p. 19).
%MAIS CITAÇÃO PRA PROCURAR
%TODO: TROCAR CITAÇÕES (USPTO, 2018, § 1207.02) para padrão ABNT (USPTO, 2018, Seção 1207.02)
%USPTO. Manual of Patent Examining Procedure. Disponível em: https://www.uspto.gov/web/offices/pac/mpep/mpep-1200.html. Acesso em: [data de acesso].

Além disso, o uso de MT em traduções de patentes representa outra aplicação comum e relativamente aceita. O Manual de Procedimentos de Exame de Patentes do Escritório de Patentes e Marcas dos Estados Unidos permite que examinadores usem MT como suporte em decisões de rejeição de patentes. \textit{"Se um documento utilizado pelo examinador para justificar uma rejeição estiver em um idioma diferente do inglês, uma tradução deve ser obtida para que o registro seja claro quanto aos fatos precisos nos quais o examinador baseia sua rejeição. [...] A tradução pode ser uma tradução automática ou um equivalente em inglês do documento não escrito em inglês"} \footnote{USPTO - United States Patent and Trademark Office. 2018. Disponível em: \url{https://mpep.uspto.gov/RDMS/MPEP/current\#/current/d0e122292.html}}.
%\cite[Seção~1207.02]{USPTO2018}.
%(USPTO, 2018, Seção 1207.02).

%"No entanto, mesmo nesse contexto, o uso de MT não é isento de riscos". No texto citado há bastante referencia sobre os riscos, entendo que é o ponto da materia evidenciar os riscos, mas não sei se devo repeti-los aqui onde tento fazer uma proposta pra aliviar o problema

\section{Avaliação da tradução automática}

Até o momento, discutimos as abordagens e implicações da tradução de máquina (MT) em diferentes contextos, incluindo o jurídico, onde os requisitos de precisão e contexto são críticos. No entanto, avaliar a qualidade e confiabilidade de uma tradução gerada automaticamente é uma tarefa complexa e multifacetada. Nesta seção, apresentamos as principais métricas de avaliação de MT e os desafios específicos associados ao uso de MT no setor jurídico.

\subsection{\textbf{Métricas quantitativas de avaliação de MT}}
Um dos métodos mais utilizados para avaliar a tradução automática envolve métricas quantitativas, que comparam a tradução gerada pelo sistema com traduções de referência. Segundo \citeonline{cheragui}, as métricas mais comumente empregadas incluem:

\begin{itemize}
\item \textbf{BLEU (\textit{BiLingual Evaluation Understudy})}: Desenvolvida por Papineni em 2001, essa métrica mede a similaridade entre a tradução da máquina e traduções de referência feitas por humanos, penalizando traduções mais curtas com o fator de penalidade de brevidade. BLEU é amplamente aceita como uma métrica objetiva definido por uma fórmula matemática que determina um fator de similaridade em n-grama \cite[p. 167]{cheragui}.

%TEM UMA FORMULA MATEMATICA AQUI MAS VOU ESCREVER ELA EM LATEX NO PROXIMO CAPITULO POIS ESSA É A MÉTRICA QUE VOU ESCOLHER
%BLEU pode ser bem similar a como você pode realmente qualificar uma NLP, já que a NLP por si só não parece ter um foco de ação em ordenação semantica/sintatica de uma frase. Até onde entendi, a NLP tentaria fazer uma predição, basicamente, tentando se aproximar de como um humano traduziria um texto. Realizar uma comparação de como o texto feito pelo computador é, e como a tradução feita pelo humano é, é evidentemente como o computador trataria Verdadeiros Positivos e falsos positivos, certo?
    
\item \textbf{WER (\textit{Word Error Rate})}: Proposta por Popovic e Ney em 2007, esta métrica calcula a distância de Levenshtein entre a tradução hipotética, feita pela máquina, e a referência, feita por um humano. Essa métrica então considera inserções, deleções e substituições necessárias para igualar as traduções. WER não permite a reorganização de palavras, limitando sua aplicabilidade em frases onde a tradução da hipótese pode ser correta, porém com ordem gramatical diferente da tradução referêncial \cite[p. 167]{cheragui}. Este tipo de rigidez quanto a flexibilidade a ordem das palavras se dá ao fato de que WER tem seu foco no reconhecimento de fala, não apenas de tradução em máquina.

%Preciso pensar aqui um pouco mais pra conseguir elaborar uma imagem, acho que ajudaria bastante poder exemplificar essas técnicas
    
%m"The cat quickly jumped over the mat."
%"The cat jumped over the mat quickly."

% A tradução tá correta, mas WER avaliaria ela penalizada erroneamente, pois a ordem das palavras está diferente 
    
\item \textbf{PER (\textit{Position-independent word Error Rate})}: A PER foi proposta anterior a WER, em 1997. Esta métrica compara as palavras feitas pela tradução de máquina, com as palavras da tradução referencial, não levando consideração a ordenação das palavras. A PER, portanto, falha em reconhecer traduções corretas onde a ordem semântica muda, mudando o sentido semântico de uma tradução \cite[p. 168]{cheragui}.

%The cat sat on a mat
%The mat sat on a cat

%PER aparentemente avaliaria bem erroneamente, pois ela identifica todas as palavras auqi, mas semanticamente tem um sentido completamente diferente
    
\item \textbf{TER (\textit{Translation Error Rate})}: Proposta subsequente a PER, em 2006, TER propõe-se a medir o número mínimo de edições para transformar a tradução da máquina em uma referência, incluindo inserção, deleção e substituição de palavras e até a movimentação de sequências de palavras. Sendo assim, uma métrica mais flexível, porém que põe em evidência vulnerabilidade quanto a manutenção contextual de uma frase. \cite[p. 168]{cheragui}.

%Imaginamos que uma tradução mude erroneamente a voz ativa de uma frase para voz passiva
%O homem atacou o idoso
%O idoso foi atacado pelo homem
%TER, não necessariamente penalizaria a frase pela ordenação das palavras, pelo menos não severamente, porém é discutivel notar que a voz ativa ou passiva tem grande relevância na hora de apresentar fatos, ainda mais relevante em contexto jurídico

\end{itemize}

\subsection{\textbf{Nuances na avaliação da eficácia da MT no contexto jurídico}}

\citeonline{VHS} apontam as dificuldades práticas na avaliação da MT em textos jurídicos. Eles destacam que as avaliações anteriores muitas vezes negligenciam a natureza dependente de contexto e subjetiva da qualidade da tradução, conforme observado por \citeonline{somers}. Somers também nota que variáveis como o comprimento das frases influenciam a contagem de erros, o que pode distorcer a avaliação da MT, já que um erro em uma frase curta pode ser proporcionalmente mais grave do que múltiplos erros em uma frase longa.
%Somers (2007)

\citeonline{wahler} sugere estabelecer requisitos mínimos de precisão como medida para reduzir riscos na utilização de MT em contextos legais, mas \citeonline{VHS} ressaltam as dificuldades práticas, dado que a precisão depende tanto da complexidade do texto original quanto do método de avaliação utilizado. Contudo, os autores advertem contra a dependência exclusiva em métricas automáticas, que, conforme \cite{callison}, são limitadas, trazendo erros comuns que poderiam ser evitados dada a intervenção feita por uma avaliação humana. \textit{"De forma mais ampla, declarações como essas desconsideram a natureza superficial das métricas automáticas e o fato de que, na prática, são medidas de similaridade que não levam em conta o efeito real ou a gravidade de quaisquer erros"} \cite{VHS}
%Wahler (2018)

%\subsubsection{}
%\paragraph{}